\documentclass[11pt]{article}

\usepackage[margin=1in]{geometry}
\usepackage[T1]{fontenc}
\usepackage[utf8]{inputenc}
\usepackage{lmodern}
\usepackage{microtype}
\usepackage{amsmath}
\usepackage{booktabs}
\usepackage{enumitem}
\usepackage[numbers,sort&compress]{natbib}
\usepackage[colorlinks=true,linkcolor=blue!50!black,citecolor=blue!50!black,urlcolor=blue!50!black]{hyperref}

\newcommand{\sys}{OpenGuide}
\newcommand{\rt}{GuidanceGrid}
\newcommand{\obc}{OBC}
\newcommand{\code}[1]{\texttt{#1}}

\title{\textbf{OpenGuide: A Control Plane for LLM Inference Trajectories}\\[2pt]
\large Declared behavior, tiered enforcement, and a replayable hash-chained record}
\author{Mahesh Vaijainthymala Krishnamoorthy\\\small\texttt{mahesh.vaikri@gmail.com}}
\date{Draft --- \today}

\begin{document}
\maketitle

\begin{abstract}
Production LLM agents are deployed with a data plane---application, orchestration,
model, inference server---but no control plane: nothing continuously compares the
behavior an operator \emph{declared} with the behavior a running trajectory
\emph{exhibits}, corrects divergence, and keeps a verifiable record of having done
so. Observability tools sense but do not act; prompt and activation-steering
methods act but do not sense against a declared setpoint or leave an auditable
record. We present \sys, a behavioral control plane that closes this loop. A
declarative, content-addressed \emph{Open Behavioral Contract} (\obc) states
invariants, phases, triggers and a recovery ladder; a compiler resolves one
contract to the strongest enforcement channel a target model permits (activation
steering, constrained decoding, scheduled re-anchoring, or structural validation)
and \emph{refuses to run} rather than silently degrade below the contract's floor;
typed sensors read behavioral state each turn; a reconciler computes drift and
dispatches recovery; and every control action is appended to a hash-chained trace
with a bit-identical replay property. We describe the design, the reference
implementation (\rt), and a conformance suite of golden vectors. We evaluate the
control path on a deterministic \emph{simulated} agent: against fixed-cadence
re-prompting at matched intervention count, reactive contract-driven re-anchoring
raises mean adherence by $+0.064$ (steady) and $+0.112$ (sparse-shock) over 50
seeds, at about $1\%$ more tokens. These are methodology-preview results from an
idealized simulator, \emph{not} evidence about real models; we state this and the
resulting threats to validity prominently, and we do not report any
capability-retention numbers because the real-model harness has not been run.
\end{abstract}

\section{Introduction}
\label{sec:intro}

Agentic LLM systems run for tens to hundreds of turns. Over such horizons,
instruction adherence erodes: constraints stated at turn zero are weighted less
as context grows, scope creeps, and tone or policy drifts. Turn-wise measurements
of this effect in multi-turn interaction are now reported in the literature
\citep{dongre2025drift}. The operational response today is largely manual:
engineers re-edit prompts, add periodic reminders, or inspect traces after the
fact.

The surrounding infrastructure is organized as a data plane. Requests flow from
application through orchestration to model and inference server, and each layer
has tooling for \emph{seeing} what happened. What is missing is the analogue of a
control plane in cluster orchestration: a component that holds a \emph{declared}
desired state, continuously measures \emph{actual} state, reconciles the two, and
records what it did. We argue that three existing families each supply one
sub-function of this loop but none closes it:

\begin{itemize}[nosep,leftmargin=*]
\item \textbf{Observability} senses, but is read-only by definition.
\item \textbf{Steering and prompting} acts---on activations
  \citep{turner2023actadd,zou2023repe,li2023iti,rimsky2023caa} or on text---but is
  typically open-loop: it emits and hopes, with no declared setpoint and no
  record that the intervention held.
\item \textbf{Governance tooling} declares policy and produces documents, but
  gates or documents rather than continuously correcting.
\end{itemize}

\sys{} composes these into one loop: \emph{declare} behavior, \emph{compile} it
to whatever enforcement the target allows, \emph{sense} behavioral state,
\emph{reconcile} drift against setpoints, \emph{actuate} a correction, and
\emph{record} every step in a tamper-evident trace.

\paragraph{Scope and honesty.}
We deliberately avoid several claims. \sys{} \emph{biases} trajectories
statistically; it does not guarantee behavior (only hard structural invariants at
the decoding and validation tiers are mechanically enforced). We do not claim a
novel control-theoretic framing of steering---control-theoretic treatments of
steering and representation editing exist
\citep{nguyen2025pidsteer,cheng2024liseco,kong2024recontrol}---nor the first time-varying
steering \citep{scalena2024dac}. Our contribution is systems-level: persistence,
addressability, cross-tier compilation, and verifiable replay of the control
record. Our quantitative results are from a simulator (Section~\ref{sec:eval}).

\paragraph{Contributions.}
\begin{enumerate}[nosep,leftmargin=*]
\item \textbf{A contract format} (\obc) that is declarative, model-independent,
  and content-addressed, so a contract hash is a deployment identity
  (Section~\ref{sec:obc}).
\item \textbf{A tiered compiler} that resolves one contract to the strongest
  available enforcement channel and refuses to run below a declared floor
  (Section~\ref{sec:compiler}).
\item \textbf{A reconciliation loop with a single write path}, in which the only
  side effect of the control channel is appending a record to a hash-chained trace
  with a bit-identical replay property (Sections~\ref{sec:reconciler}
  and~\ref{sec:trace}), together with a cross-implementation conformance suite.
\item \textbf{A simulated evaluation protocol with an explicit kill criterion},
  and an honest account of what it does and does not show
  (Section~\ref{sec:eval}).
\end{enumerate}

\section{Related work}
\label{sec:related}

\paragraph{Activation steering.}
Adding vectors to the residual stream to shift behavior is established
\citep{turner2023actadd,rimsky2023caa,li2023iti}, with a broader
representation-engineering view of reading and controlling internal state
\citep{zou2023repe}. Conditional activation steering applies an intervention only
when a detector fires \citep{lee2024cast}; dynamic activation composition varies
steering strength across generation steps \citep{scalena2024dac}; and a
feedback-controller formulation treats steering as closed-loop control over
activations \citep{nguyen2025pidsteer}. Linear semantic control
\citep{cheng2024liseco} and a control-perspective treatment of representation
editing \citep{kong2024recontrol} similarly bring control theory to generation.
Selective Steering studies norm-preserving
control with discriminative layer selection \citep{dang2026selective}. \sys{} uses
norm-preserving injection of this kind at its open-weight tier, but its focus is
different: it is the \emph{surrounding system}---contract, compiler, state,
trace---not a new steering rule.

\paragraph{Constrained decoding.}
Grammar- and automaton-guided generation provides hard structural guarantees on
output \citep{willard2023outlines}. \sys{} treats this as its decoding tier and as
the home of \emph{hard} invariants.

\paragraph{Multi-turn drift.}
\citet{dongre2025drift} formalize turn-wise divergence in multi-turn interaction.
We adopt the idea of summarizing a run by cumulative deficit below a floor
(Section~\ref{sec:eval}); we do not claim to reproduce their methodology.

\paragraph{Where \sys{} sits.}
Relative to observability, \sys{} adds actuation. Relative to steering and
prompting, it adds a declared setpoint, sensing, and a verifiable record.
Relative to governance, it adds continuous correction. Removing any one leg
collapses it into an existing category: without actuation it is an observability
tool, without sensing a guardrail, without the contract a steering library.

\section{Design}
\label{sec:design}

\subsection{Open Behavioral Contracts}
\label{sec:obc}

An \obc{} is a TOML document with sections for identity (persona, boundaries),
\emph{invariants}, \emph{phases}, \emph{triggers}, a \emph{recovery} ladder, and a
signal table that binds abstract signal names to tier-specific artifacts. Each
invariant names a \emph{measure} whose prefix types its honesty:
\code{structural:} (deterministic check), \code{probe:} (linear probe on the
residual stream, open weights only), \code{external:} (webhook or registered
callable), or \code{judge:} (LLM-as-judge, explicitly tagged stochastic). An
invariant carries either a soft \code{floor} that contributes to drift or
\code{hard = true}, meaning it is enforced and never advisory.

A contract is content-addressed: $\text{hash}(c)=\mathrm{sha256}(\mathrm{canon}(c))$
where $\mathrm{canon}$ sorts keys, rounds floats to nine significant digits, and
emits UTF-8 JSON without insignificant whitespace. The same contract therefore
hashes identically across languages and TOML libraries; the hash is the
deployment identity (the proxy is activated by presenting it in a request
header).

\subsection{Enforcement tiers and the capability resolver}
\label{sec:compiler}

\begin{table}[t]
\centering
\caption{Enforcement tiers. Higher fidelity requires more access to the model.}
\label{tab:tiers}
\begin{tabular}{@{}llll@{}}
\toprule
Tier & Channel & Requires & Nature \\
\midrule
T0 & Activation steering & open weights & soft, statistical \\
T1 & Constrained decoding & local decoding control & hard (structural) \\
T2 & Scheduled re-anchoring & prompt access (closed APIs) & soft, text-level \\
T3 & Structural validation & output access (universal) & hard (detect/retry) \\
\bottomrule
\end{tabular}
\end{table}

Each target model has a static \emph{capability descriptor} listing the tiers it
supports (e.g.\ closed APIs: T2, T3; vLLM and HF Transformers: T0--T3; llama.cpp:
T1--T3). The resolver selects the strongest supported tier that is at least as
strong as the contract's \code{enforcement\_floor}. If none exists it raises an
error: \emph{the runtime refuses to run rather than silently degrading
enforcement.} Unknown providers are assumed to be closed-API-equivalent (T2/T3).
The compiler then emits tier artifacts: T2 prompt fragments and a re-anchor plan,
T3 validators, probe specifications, a trigger table, and a recovery plan.

\subsection{The gain grid (open-weight tier)}
\label{sec:grid}

At T0, steering strength is a \emph{total} scalar field $G[\ell,t]$ over the
(layer $\times$ step) lattice: every coordinate is defined, unset cells are zero.
Storage is sparse globally and dense locally (fixed-width time frames holding
dense layer-by-frame tiles). Fields are closed under addition, scalar scaling,
masking, shifting and cropping, so grids \emph{compose}: a tenant temperament and
a task schedule combine as $G_3=G_1\oplus\alpha G_2$ and the result is itself a
hashable artifact. Injection is norm-preserving,
\[
 h' = \frac{\lVert h\rVert}{\lVert h+g\,v\rVert}\,(h+g\,v),
\]
and a per-layer cap $\Sigma$ on summed absolute gain across signals is enforced
\emph{mechanically} (columns are scaled down and the correction recorded), not
advisorily. Steering vectors live in a content-addressed \emph{signal pack}
referenced by hash only.

\subsection{Sensing and the reconciler}
\label{sec:reconciler}

Each turn the reconciler: (1) reads every invariant's measure, yielding a value
or a \emph{gap}; (2) computes drift; (3) runs hard validators; (4) evaluates
state-predicate triggers; (5) dispatches recovery from the ladder; (6) assembles
the T2 re-anchor text to inject on the next turn; and (7) appends one behavioral
state record to the trace.

Drift is the weighted \emph{sum} of soft-floor violations, clipped to $[0,1]$:
\[
 d_t=\min\!\Big(1,\ \sum_{i\in\mathcal S} w_i\,\max(0,\,f_i-r_{i,t})\Big),
\]
where $f_i$ is the floor, $w_i$ the invariant's drift weight and $r_{i,t}$ its
reading; invariants without a positive floor and invariants whose reading is a gap
do not contribute. A sum (rather than a mean) is used so that a healthy invariant
contributes zero and never dilutes another's violation. The recovery ladder maps
$d_t$ to an action (\code{re\_anchor}, \code{escalate}, \code{halt}); a hard
validation failure forces \code{re\_anchor}.

Two engineering rules matter. First, \emph{the data plane never blocks on the
control plane}: a sensor timeout or failure yields a recorded \emph{sensor gap}
and the run proceeds. Second, measures are \emph{typed by honesty}: each declares
whether it is deterministic or stochastic and a latency class, and a judge-based
measure never silently substitutes for a structural one.

\subsection{The trace: hash-chained, replayable}
\label{sec:trace}

A trace is a \emph{manifest}---model, contract hash and body, pack hash, seed,
runtime version---plus an append-only record list. With $\mathrm{chash}$ the
canonical hash,
\[
 \text{head}_0=\mathrm{chash}(\text{manifest}),\qquad
 \text{head}_i=\mathrm{chash}\big(\{\text{prev}:\text{head}_{i-1},\ \dots\text{event}_i\}\big).
\]
Verification recomputes the chain and compares the final head. Editing,
reordering, inserting or deleting any record changes the head. \emph{Replay
identity}: the same manifest and the same recorded events yield a bit-identical
head. The central invariant, enforced by construction, is that run state mutates
only through the control channel and the channel's only side effect is appending a
record---there is no side door.

\paragraph{What the chain does and does not give.}
The chain provides \emph{tamper-evidence relative to a trusted head}. It does not
provide authenticity: a party with write access to the whole trace can recompute a
self-consistent chain. Anchoring the head externally (signing, a transparency log,
or write-once object storage) is required for non-repudiation and is outside the
current implementation.

\section{Implementation}
\label{sec:impl}

\rt{} is a Python reference implementation released under the Apache-2.0 license. Packages cover the
core (canonical encoding, gain field, grid, trace), the \obc{} loader and schema
validator, the compiler and capability resolver, the four channels, typed
sensors, the reconciler, and a SQLite content-addressed store with a fleet index
(runs queryable by contract hash, drift and halt status).

A FastAPI proxy exposes an OpenAI-compatible \code{/v1/chat/completions}
endpoint (including streaming). It is pass-through by default; presenting a
contract hash in a request header activates governance. The proxy reassembles
content and tool-call deltas from the stream, maps a conversation identifier to a
persistent reconciler and trace so that state carries across calls, injects the
pending re-anchor on the next turn, and emits one OpenTelemetry span per turn with
\code{gen\_ai.behavior.*} attributes. Client response bytes are not rewritten. A
tool-policy layer can additionally allow, or require approval for, named tool
calls, recording the decision and reason (but not arguments) as a trace event.

The T0 channel registers PyTorch forward hooks on decoder layers; a single forward
pass is one step, so it composes with \code{generate()}. The T1 channel implements
logit masking with allowlist and sequence constraints, and names a grammar engine
\citep{willard2023outlines} as the intended production backend (not installed in
the reference environment).

\paragraph{Conformance.}
Because hashes are only useful if independent implementations agree, the
specification ships golden vectors: ten canonical-encoding vectors, five
chain-verification vectors (valid, and tampered by edit, reorder and delete), and
contract-hash vectors. An independent runner that uses only the written
specification re-derives every hash; the suite currently passes 28 of 28 checks,
and CI regenerates the vectors from the reference implementation and fails on
drift between specification and implementation. The runtime test suite (305 tests)
passes in our environment.

\section{Evaluation}
\label{sec:eval}

\paragraph{Status.}
\textbf{All results in this section are simulated.} They come from a deterministic
mock agent, not a language model. They fix the methodology and exercise the real
control path end to end; they do not measure real-model behavior, and their
magnitudes should not be read as expected real-world gains.

\subsection{Setup}
A coding-agent adherence signal $a_t$ evolves over $T=60$ turns under a
disturbance: slow, late-growing decay plus irregular scope-creep shocks. Two
regimes are used: \emph{steady} (continuous decay, regular shocks) and
\emph{sparse} (mostly healthy, rare large shocks). An intervention restores a
fixed amount of adherence on the following turn and costs a fixed number of
tokens; base cost is 180 tokens/turn and each intervention adds 130. Three arms
share these constants, so only \emph{timing} differs:
\emph{none}; \emph{naive}, which re-prompts at a fixed cadence; and
\emph{contract}, which runs the actual \rt{} reconciler and re-anchors only when
drift crosses the contract threshold. The naive cadence is set to
$\mathrm{round}(T/n_c)$, where $n_c$ is the contract arm's intervention count,
matching the number of interventions.

\paragraph{Kill criterion (declared in advance).}
If naive fixed-cadence re-prompting matches contract-driven re-anchoring at equal
cost, \sys{} should be scoped to observability and the open-weight tiers, and say
so.

\subsection{Intervention benchmark}

\begin{table}[t]
\centering
\caption{Single draw (seed 7), simulated. Adherence floor $0.85$.}
\label{tab:t1}
\begin{tabular}{@{}llrrr@{}}
\toprule
Regime & Arm & Mean adherence & \% turns in violation & Tokens \\
\midrule
steady & none     & 0.303 & 88\% & 10{,}800 \\
steady & naive    & 0.819 & 57\% & 11{,}710 \\
steady & contract & \textbf{0.869} & \textbf{38\%} & 11{,}840 \\
sparse & none     & 0.483 & 73\% & 10{,}800 \\
sparse & naive    & 0.802 & 62\% & 11{,}190 \\
sparse & contract & \textbf{0.913} & \textbf{17\%} & 11{,}320 \\
\bottomrule
\end{tabular}
\end{table}

Table~\ref{tab:t1} shows one draw. To avoid resting on a single seed we repeat the
identical protocol over 50 seeds (Table~\ref{tab:sweep}); the sweep script is in
the repository.

\begin{table}[t]
\centering
\caption{50-seed sweep, simulated. $\Delta$ is contract minus naive.}
\label{tab:sweep}
\begin{tabular}{@{}lrrrrr@{}}
\toprule
Regime & Contract adh. & Naive adh. & $\Delta$ mean (sd) & Win rate & $\Delta$ tokens \\
\midrule
steady & 0.877 & 0.813 & $+0.064$ ($0.032$) & 100\% & $+130$ \\
sparse & 0.884 & 0.773 & $+0.112$ ($0.039$) & 100\% & $+130$ \\
\bottomrule
\end{tabular}
\end{table}

In both regimes the contract arm holds higher adherence, and the gap is larger
when disturbances are rare and large: a fixed cadence either spends budget
re-prompting a healthy agent or misses a shock, whereas reactive re-anchoring
spends only when drift demands it. Note that the contract arm spends 130 more
tokens per run (one intervention's worth, about $1.1\%$), so ``equal cost'' holds
to that tolerance, not exactly.

\paragraph{Sensitivity: a poorly tuned contract loses.}
The result depends on the contract's \code{drift\_threshold}. The benchmark
contract uses a floor-targeting threshold of $0.05$. With the default
coding-agent threshold of $0.27$ the reconciler re-anchors only near catastrophe,
and fixed-cadence re-prompting holds the floor better. The value is therefore in
\emph{reactive timing}, and threshold selection is a first-class part of
authoring a contract. This also means the headline comparison involves a
benchmark-specific tuning choice that favors the contract arm.

\subsection{Drift benchmark}
Following the idea of summarizing a run by deficit below the floor, we report
area-under-drift (cumulative floor deficit), peak drift, and turns in violation
(Table~\ref{tab:drift}, seed 7). The governed trace verifies in both regimes, so
every recovery is recorded and replayable.

\begin{table}[t]
\centering
\caption{Governed vs.\ ungoverned, simulated, seed 7.}
\label{tab:drift}
\begin{tabular}{@{}llrrrr@{}}
\toprule
Regime & Arm & Area-under-drift & Peak drift & Turns in violation & Recovery (turns) \\
\midrule
steady & ungoverned & 33.19 & 0.85 & 53 & never \\
steady & governed   & 1.19  & 0.18 & 23 & 1 \\
sparse & ungoverned & 23.13 & 0.81 & 44 & never \\
sparse & governed   & 0.86  & 0.26 & 10 & 4 \\
\bottomrule
\end{tabular}
\end{table}

\subsection{Overhead}
In-process, against a mock upstream, governance adds about $0.25$\,ms at the
median (pass-through $1.64$\,ms, governed $1.89$\,ms; $n=200$), against a design
gate of $25$\,ms. This excludes network and real-model latency and says nothing
about judge-sensor cost, which dominates in practice.

\subsection{Integration demonstrations on public data}
Two reference implementations route a \emph{scripted} agent through the proxy
using live public data---SEC EDGAR company facts, and the CISA Known Exploited
Vulnerabilities catalog---and demonstrate action governance: a low-risk tool call
(\code{crm.create\_account\_note}, \code{siem.create\_case}) is allowed and
recorded, while a consequential one (\code{erp.update\_vendor\_risk},
\code{edr.isolate\_host}) is blocked as approval-required. The run trace verifies
and contains the state and tool-policy events. Because the agents are scripted,
these demonstrate integration mechanics and auditability, \emph{not} model
quality.

\section{Limitations and threats to validity}
\label{sec:limits}

\begin{itemize}[nosep,leftmargin=*]
\item \textbf{The evaluation is simulated and favorable by construction.} The
  restore model is idealized (an intervention restores a fixed amount of
  adherence), the disturbance is synthetic, and the adherence signal is observed
  without noise. Real models may respond to re-anchoring non-monotonically, and
  real judge sensors are noisy and costly. The \emph{shape} of the result is
  plausible; its magnitude is not evidence.
\item \textbf{Tuning asymmetry.} The contract arm uses a threshold chosen to hold
  the benchmark floor, while the naive arm's only free parameter is cadence. A
  different tuning could reverse the outcome, as the sensitivity analysis shows.
\item \textbf{No capability-retention results.} Whether T0 steering at deployed
  gains preserves benchmark capability (e.g.\ MMLU, GSM8K) is unmeasured. The pack
  evaluation cards have empty slots for these numbers and the real-model gate is
  not implemented. We make no retention claim.
\item \textbf{Statistical, not guaranteed.} Soft floors are reconciled
  statistically; only hard invariants at T1/T3 are mechanical.
\item \textbf{T2 is text-level.} On closed APIs, enforcement is prompt injection
  and post-hoc validation, with the corresponding limits.
\item \textbf{Trace authenticity.} See Section~\ref{sec:trace}: the chain is
  tamper-evident only against a trusted head.
\item \textbf{Sensor fidelity.} Drift is only as good as the measures behind it;
  stochastic judge measures add variance, and sensor gaps are recorded, not
  repaired.
\item \textbf{Single reference implementation.} Conformance vectors cover
  canonical encoding, chain verification and contract hashing, but a second
  independent runtime has not been built.
\end{itemize}

\section{Conclusion and future work}
\label{sec:conclusion}

\sys{} frames reliable agent behavior as a control problem: declare it, enforce it
at the strongest available tier, sense it, reconcile, and keep a record that can
be verified and replayed. The contract format, the refuse-to-degrade compiler and
the single-write-path trace are the parts we consider durable; the quantitative
claim is, for now, a simulated methodology preview.

The decisive next step is to replace Section~\ref{sec:eval} with real-model
results: the intervention benchmark on a closed and an open model at matched token
budgets, and the capability-retention ladder on an open-weight model. Per our
kill criterion these will be reported regardless of outcome. Further work includes
externally anchored (signed) trace heads, a second independent runtime, a
learned re-anchor schedule evaluated against constant baselines, and a
standardized \code{gen\_ai.behavior.*} attribute set for telemetry.

\section*{Reproducibility}
The specification, reference implementation, conformance vectors and benchmark
scripts are in the \sys{} repository. From the repository root:
\code{make bench} (single-seed tables), \code{python benchmarks/seed\_sweep.py}
(50-seed sweep), \code{python spec/conformance/run\_conformance.py}
(28 checks), and \code{python scripts/bench\_proxy.py} (overhead).

\bibliographystyle{plainnat}
\bibliography{references}

\end{document}
